\section{Machine-checked verification in Lean 4}
\label{sec:formal}

Each rung of this paper corresponds to a named Lean~4 declaration of a public
development~\cite{LeanRepo}. The development consists of three independent trees: the analysis core and the
exact-criterion rung live in \lean{wC}, the Jensen rung in \lean{wF}, and the weak-constant
$8/\mathrm{e}$ variant in \lean{wM} (paths below are relative to \lean{lean_sharp/}). In each tree the model
-- the rotations, the evolution, \lean{Admissible}, and the cost $T=\texttt{cost}\,\theta$ --
is fixed in \lean{RobustZ/Statement.lean}, and the scalar flatness~\eqref{eq:hflat} is derived
from the matrix flatness (\lean{h_flat}) rather than assumed. Write $q=2N+2$ and
$\tauvar=T/2$.

\subsection{The machine-checked statements}

Each rung corresponds to a declaration of the development, stated here with its exact
hypotheses so that it can be compared with Table~\ref{tab:formal-correspondence} row by row.
The elementary rung is \lean{elementary_bound}, with \lean{elementary_bound_pi} for the case
$\phi=\pi$: for every $0<\phi\le\pi$ and every order-$N$ admissible sequence,
$T\ge2\,(2\,q!\,\sin^{2}(\phi/4))^{1/q}$, and $T\ge2\,(q!)^{1/q}$ at $\phi=\pi$. The Jensen
rung is \lean{jensen_rung_final}, with \lean{jensen_rung_pi_final} at $\phi=\pi$: with
$s=2\sin^2(\phi/4)$, $T\ge\pi q/[\mathrm{e}\,(2/s)^{1/q}]$, and \lean{jensen_ratio_tendsto}
records that this per-order bound divided by $N$ tends to $2\pi/\mathrm{e}$. The exact
criterion is \lean{new_rung_final}: for $N\ge1$, $0<\phi\le\pi$ and $1\le\tauvar<q=2N+2$, the
inequality $\mathrm{nrCcal}(\tauvar,q)\,q^2\,\mathrm{nrE2star}(\tauvar,q)<\sin^2(\phi/4)$,
with the closed forms~\eqref{eq:besseldefs}--\eqref{eq:ccal}, implies
$2\tauvar<\mathrm{cost}\,\theta$. The certified instances \lean{cert_N2}--\lean{cert_N4} and
\lean{cert_N5_tight}--\lean{cert_N12_tight} (Table~\ref{tab:certs}) and the closed form
$T\ge4(N+1)(1-\mathrm{nrU}(q))$ (\lean{nr_closed_form}, i.e.~\eqref{eq:closedform}) are proved
from it with nothing else added.

\subsection{Statement-by-statement correspondence}

Table~\ref{tab:formal-correspondence} lists the correspondence between the objects of this
paper and the declarations that check them.

\begin{table}[t]\centering\footnotesize
\caption{Correspondence between the objects of this paper and the declarations that check them.
File paths are relative to \lean{lean_sharp/}; declaration names are exact.}
\label{tab:formal-correspondence}
\begin{tabular}{@{}p{0.30\linewidth}p{0.36\linewidth}p{0.30\linewidth}@{}}
\toprule
\textbf{Paper object} & \textbf{Lean declaration} & \textbf{File}\\
\midrule
\raggedcell Model, \lean{Admissible} and \lean{cost} & \raggedcell \lean{Admissible}, \lean{cost} & \raggedcell \lean{wC/RobustZ/Statement.lean} (same in \lean{wF}, \lean{wM})\\
\raggedcell Direct bound (8), Theorem~\ref{thm:direct} & \raggedcell paper-level proof (in text) & \raggedcell Sec.~\ref{sec:elementary}\\
\raggedcell Carrier flatness (20), moments (39) & \raggedcell \lean{h_flat}, \lean{h_moments} & \raggedcell \lean{wC/RobustZ/Flatness.lean}, \lean{wC/RobustZ/ExpSum.lean}\\
\raggedcell Elementary bound, Theorem~\ref{thm:elementary} & \raggedcell \lean{elementary_bound}, \lean{elementary_bound_pi} & \raggedcell \lean{wC/RobustZ/ElementaryBound.lean}\\
\raggedcell Jensen bound, Theorem~\ref{thm:jensen}, with its analytic core & \raggedcell \lean{jensen_rung_final}, \lean{jensen_rung_pi_final}, \lean{jensen_ratio_tendsto} & \raggedcell \lean{wF/RobustZ/JensenBridge.lean}, \lean{wF/RobustZ/JensenRung.lean}\\
\raggedcell Coefficient bound L1, tail sum L2 & \raggedcell \lean{fcoef_sharp}, \lean{bessel_L1}, \lean{bessel_L2} & \raggedcell \lean{wC/RobustZ/BesselTail.lean}\\
\raggedcell Realization L3, Fourier inversion & \raggedcell \lean{kap_L1_le}, \lean{realize_L1}, \lean{inversion_on_band} & \raggedcell \lean{wC/RobustZ/RealizeLemma.lean}, \lean{wC/RobustZ/InversionId.lean}\\
\raggedcell Chebyshev derivative bundle D1 & \raggedcell \lean{cb_term1_le}, \lean{cb_term2_le}, \lean{cb_hasSum0}--\lean{cb_hasSum2} & \raggedcell \lean{wC/RobustZ/ChebBundle.lean}\\
\raggedcell Exact criterion L4/L5 and closed form, Theorem~\ref{thm:newrung} & \raggedcell \lean{nr_sup_le}, \lean{nr_D1_le}, \lean{nr_D2_le}, \lean{nr_kappa_le}, \lean{nr_main_rung}, \lean{nr_closed_form} & \raggedcell \lean{wC/RobustZ/NewRung.lean}\\
\raggedcell Wrapper and finite-$N$ certificates & \raggedcell \lean{new_rung_final}, \lean{cert_N2}--\lean{cert_N4}, \lean{cert_N5_tight}--\lean{cert_N12_tight} & \raggedcell \lean{wC/RobustZ/RungFinal.lean}, \lean{wC/RobustZ/RungTight.lean}\\
\raggedcell Weak-constant $8/\mathrm{e}$ variant & \raggedcell \lean{eight_over_e_bound} & \raggedcell \lean{wM/RobustZ/EightOverE.lean}\\
\bottomrule
\end{tabular}
\end{table}

\subsection{Scope of the verification}

Machine-checked: the model and the carrier flatness; the elementary bound with its exact
constant $2$, at every order and every $0<\phi\le\pi$; the Jensen bound with its exact
constant $\pi/\mathrm{e}$ at every order and every $0<\phi\le\pi$; the exact criterion with its
closed forms, hence every entry of Table~\ref{tab:certs}; the $8/\mathrm{e}$ variant with
$C=10^8$; and the passage to the coefficient, as follows. \lean{c_ge_four} proves
$4\le\liminf_{N}T_{\min}(N,\phi)/N$ under two side conditions, that an admissible sequence
exists at every order and that $T_{\min}(N,\phi)/N$ is eventually bounded above, and the
dichotomy \lean{c_ge_four_or_grows} keeps only the first and concludes either that same bound
or that $T_{\min}(N,\phi)/N$ diverges to $+\infty$. The closed form~\eqref{eq:closedform} is
\lean{nr_closed_form}, whose hypothesis is the existence of truncation data at the bandwidth
that form prescribes (constructed by \lean{nrCheb_of_truncation}), rather than an
unconditional per-sequence statement. Not machine-checked: the numerical material of
Sec.~\ref{sec:numerics}. The low-order sequences found by local search and their
re-verified costs are upper bounds computed outside Lean; they enter no proof, and no
optimality or small-$N$ statement about them is formalized. The audit ran on the development sources: a source scan finds no \lean{sorry},
\lean{admit}, or \lean{axiom} in the project files, and the headline theorems depend only
on the three standard logical axioms \lean{propext}, \lean{Classical.choice} and
\lean{Quot.sound} (see the \lean{Audit*.lean} files, e.g.\
\lean{wC/RobustZ/AuditRungFinal.lean}). Statements, build logs and the raw
\texttt{\#print axioms} output are collected at the project page, and the repository with the
sources is cited as~\cite{LeanRepo}.

\subsection{Code, data, and reproducibility}
\label{sec:resources}

Everything used in this paper is public: the proofs are a Lean development, the numerical sequences
and their independent re-verification are a data set with a browser front end, and the figures and
tables of Sec.~\ref{sec:numerics} are regenerated from that data by short scripts.
Table~\ref{tab:resources} collects the four entry points.

\begin{table}[t]\centering\footnotesize
\caption{Public resources for this paper. The Lean development is self-contained: a fresh clone
reproduces its statements with the three commands of its \texttt{README}. The numerical
re-verification is likewise reproducible from the last entry, which carries the sequences at the
precision used in Table~\ref{tab:num}. All addresses were accessed on 8 October 2026.}
\label{tab:resources}
\begin{tabular}{@{}p{0.22\linewidth}p{0.68\linewidth}@{}}
\toprule
\textbf{Resource} & \textbf{Content} \\
\midrule
\raggedcell \url{https://show.dytchem.cn/bloch/} &
\raggedcell Bloch-sphere animation and the full angle lists of every sequence of
Table~\ref{tab:num}; browser-side re-computation of the nominal mismatch, of the deviation
$\delta(\eps)$ and of its slope, with the reference sequence shown side by side.~\cite{DataRepo} \\
\raggedcell \url{https://show.dytchem.cn/lean/} &
\raggedcell Verification status of the formal development: the exact Lean statements of the
bounds, the source scan for \lean{sorry} and \lean{axiom}, and the raw \lean{lake build} and
\texttt{\#print axioms} output for the headline theorems.~\cite{LeanRepo} \\
\raggedcell \href{https://github.com/Dytchem/robust-z-lower-bound-lean}{GitHub repository} &
\raggedcell Lean~4 with Mathlib, pinned by \lean{lean-toolchain} and
\lean{lake-manifest.json}; \lean{VERIFICATION.md} records two from-scratch verifications on
three trees: \lean{lean_sharp/wC} carries the analysis core, the exact-criterion rung and its
certificates (54 modules); \lean{lean_sharp/wF} carries the Jensen rung and its bridge
(35 modules); \lean{lean_sharp/wM} carries the $8/\mathrm{e}$ variant (47 modules); all build
green with no \lean{sorry} or \lean{axiom} in the sources.~\cite{LeanRepo} \\
\raggedcell \href{https://show.dytchem.cn/files/code.zip}{Verification scripts and data} &
\raggedcell \lean{verify_num.py}, which rebuilds every propagator with
\lean{scipy.linalg.expm} and re-evaluates the Taylor coefficients of Sec.~\ref{sec:numerics}
at $50$ decimal digits; \lean{make_outputs.py}, which regenerates Figs.~\ref{fig:num}
and~\ref{fig:phys} and the printable angle lists of Appendix~\ref{app:sequences} from the
verified data; the data themselves (\lean{num_results.json}, \lean{solutions_largeN.json}); and
a \lean{README}.~\cite{DataRepo} \\
\bottomrule
\end{tabular}
\end{table}

Each tree builds green from its pinned toolchain and manifest alone.
The angle lists of Appendix~\ref{app:sequences} are stored in the convention in which they were
verified, rounded to $12$ decimals of $\pi$; the rounding is checked, not assumed: feeding those
rounded lists back into~\eqref{eq:rotations}--\eqref{eq:robustness} again gives a nominal
mismatch below $9\times10^{-10}$ and a flatness residual below $1.5\times10^{-10}$, which are the
two accuracy statements of Table~\ref{tab:num}.
